\begin{center}
\textbf{Abstract}
\end{center}

Aquest Treball de Final de Grau presenta el disseny i la implementació d'un intèrpret educatiu pas a pas per a un subconjunt del llenguatge de programació funcional Haskell, desenvolupat íntegrament en Haskell. El projecte neix amb l'objectiu de resoldre la desconnexió cognitiva que pateixen els estudiants en introduir-se al paradigma funcional pur, on l'avaluació mandrosa (lazy evaluation) i les opacitats dels compiladors tradicionals, com GHCi, dificulten l'aprenentatge. L'arquitectura del sistema es divideix en tres grans mòduls: un analitzador lèxic i sintàctic construït des de zero mitjançant combinadors de parsers monàdics, evitant l'ús de generadors de codi externs; un motor d'inferència de tipus estàtic basat en l'algorisme de Hindley-Milner i la unificació de Robinson; i un avaluador mandrós fonamentat en la reducció de grafs i l'ús de supercombinadors. Els resultats demostren la viabilitat d'una implementació modular i declarativa capaç de processar estructures complexes com l'aplicació parcial, la currificació i les clausures lèxiques. Aquesta eina assoleix una fita tècnica de gran rigor i esdevé un recurs pedagògic clau per fer transparent el procés d'execució i facilitar la comprensió de les bases teòriques de la programació funcional.


This Bachelor's Thesis presents the design and implementation of an educational, step-by-step interpreter for a subset of the functional programming language Haskell, developed entirely in Haskell. The project stems from the objective of bridging the cognitive gap students face when introduced to the pure functional paradigm, where lazy evaluation and the black-box nature of traditional compilers, such as GHCi, hinder the learning process. The system's architecture is divided into three main modules: a lexical and syntactic analyzer built from scratch using monadic parser combinators, without relying on external code generators; a static type inference engine based on the Hindley-Milner algorithm and Robinson's unification; and a lazy evaluator founded on graph reduction and supercombinators. The results demonstrate the viability of a modular and declarative implementation capable of processing complex structures such as partial application, currying, and lexical closures. This tool achieves a highly rigorous technical milestone and serves as a key pedagogical resource to make the execution process transparent and facilitate the understanding of the theoretical foundations of functional programming.

Este Trabajo de Final de Grado presenta el diseño y la implementación de un intérprete educativo paso a paso para un subconjunto del lenguaje de programación funcional Haskell, desarrollado íntegramente en Haskell. El proyecto nace con el objetivo de resolver la desconexión cognitiva que sufren los estudiantes al introducirse en el paradigma funcional puro, donde la evaluación perezosa (lazy evaluation) y la opacidad de los compiladores tradicionales, como GHCi, dificultan el aprendizaje. La arquitectura del sistema se divide en tres grandes módulos: un analizador léxico y sintáctico construido desde cero mediante combinadores de parsers monádicos, sin dependencias de generadores de código externos; un motor de inferencia de tipos estático basado en el algoritmo de Hindley-Milner y la unificación de Robinson; y un evaluador perezoso fundamentado en la reducción de grafos y el uso de supercombinadores. Los resultados demuestran la viabilidad de una implementación modular y declarativa capaz de procesar estructuras complejas como la aplicación parcial, la currificación y las clausuras léxicas. Esta herramienta alcanza un hito técnico de gran rigor y se convierte en un recurso pedagógico clave para hacer transparente el proceso de ejecución y facilitar la comprensión de las bases teóricas de la programación funcional.